\section{Results}
\label{sec:results}

\subsection{Window size optimization identifies 8,192\,\bp{} as optimal}
\label{sec:results_window}

To determine the scoring window that best captures variant context, we evaluated
four window sizes (1,024, 2,048, 4,096, and 8,192\,\bp) on ClinVar-annotated
variants in three well-characterised genes (BRCA1, TP53, CHEK2). The
8,192\,\bp{} window achieved the highest mean AUROC (0.9921) for discriminating
pathogenic from benign variants and was selected for all subsequent analyses
(\Cref{fig:supp_ablation}). Performance improved monotonically with window size
across all three genes, suggesting that flanking-sequence context---including
intronic and regulatory regions---contributes to variant effect prediction.

\subsection{DMS calibration confirms Evo2 captures functional impact}
\label{sec:results_dms}

We calibrated Evo2 $\Delta$LL scores (\Cref{eq:delta_ll}) against deep
mutational scanning data for four control genes (\Cref{fig:dms}A--D). Evo2 showed significant Spearman
correlations with experimentally measured functional scores in all four genes:
BRCA1 ($\rho = +0.515$, $p < 10^{-261}$, $n = 3{,}884$),
DNMT3A ($\rho = +0.457$, $p < 10^{-39}$, $n = 746$),
CHEK2 ($\rho = +0.310$, $p < 10^{-108}$, $n = 4{,}882$),
and TP53 ($\rho = -0.211$, $p < 10^{-29}$, $n = 2{,}827$). The negative
TP53 correlation is expected: the Giacomelli nutlin-3 assay assigns high scores
to loss-of-function variants, inverting the sign convention.

Multi-tool comparison on the same variant sets revealed that Evo2's $|\rho|$
was competitive with supervised tools (\Cref{fig:dms}E; Supplementary Table~S3).
For BRCA1, Evo2 ($|\rho| = 0.515$) was comparable to AlphaMissense (0.562) and
exceeded CADD (0.499) and REVEL (0.470). AlphaMissense achieved the highest
$|\rho|$ for DNMT3A (0.672 vs.\ Evo2 0.457) and TP53 (0.414 vs.\ 0.211),
consistent with its specialisation for missense variants. SpliceAI and ncER
showed the weakest DMS correlations ($|\rho| < 0.28$), as expected given that
DMS variants are predominantly coding.

\subsection{ClinVar benchmarking demonstrates clinical-grade classification}
\label{sec:results_clinvar}

We benchmarked Evo2 against CADD, AlphaMissense, and REVEL on ClinVar variants
with $\geq$2-star review status across six genes with sufficient pathogenic and
benign annotations (\Cref{fig:clinvar}; Supplementary Table~S5).

\subsubsection*{Per-tool AUROCs.}

On its full scoreable variant set, Evo2 achieved AUROCs of 0.909--0.9996
(\Cref{fig:clinvar}A,B). Evo2 was the highest-performing tool for two genes:
CHEK2 (0.9996, 95\% CI [0.999, 1.000]; vs.\ CADD 0.9986) and TP53 (0.9887
[0.983, 0.993]; vs.\ CADD 0.9882, AlphaMissense 0.988). For ATM (0.9955),
BRCA1 (0.9879), and DNMT3A (0.9962), Evo2 was slightly below CADD
(0.9989, 0.9951, 1.000 respectively) but within overlapping confidence
intervals. TERT showed the lowest Evo2 AUROC (0.9089 [0.832, 0.974]), consistent
with this gene's non-coding nature: many TERT ClinVar variants are promoter
mutations, including gain-of-function hotspots (C228T, C250T) that violate the
loss-of-function assumption underlying $\Delta$LL scoring.

\subsubsection*{Matched-variant analysis.}

On the intersection of variants scored by all four tools (Evo2 $\cap$ CADD
$\cap$ AlphaMissense $\cap$ REVEL), sample sizes were substantially reduced
($n = 18$--$335$; \Cref{fig:clinvar}C; Supplementary Table~S5). These matched
sets were dominated by missense variants, which inherently favours
AlphaMissense and REVEL. On matched sets, AlphaMissense led for TP53
(0.988 vs.\ Evo2 0.894; DeLong $p = 2.0 \times 10^{-5}$) and BRCA1 (0.960 vs.\
0.932; $p = 0.042$). Importantly, the DeLong comparison between CADD and Evo2
was not significant for any gene ($p > 0.08$), indicating that Evo2's zero-shot
approach achieves comparable discrimination to CADD's supervised framework on
shared variant sets.

\subsection{Orthogonal signal enables ensemble improvement}
\label{sec:results_orthogonality}

Evo2 scores showed moderate to low correlation with supervised tools, indicating
complementary information content (\Cref{fig:orthogonality}A). Spearman
correlations between Evo2 and CADD ranged from 0.414 (TERT) to 0.828 (TP53);
between Evo2 and SpliceAI from 0.053 (DNMT3A) to 0.446 (CHEK2). This
orthogonality motivated an ensemble approach.

Five-fold cross-validated logistic regression combining Evo2 and CADD scores
as features consistently matched or exceeded the best single tool
(\Cref{fig:orthogonality}B; Supplementary Table~S5). The largest gain
was observed for TP53, where the ensemble AUROC (0.9922) exceeded both Evo2-only
(0.9882) and CADD-only (0.9885) by 0.38\%. For CHEK2, the ensemble (0.9998)
improved on an already near-perfect Evo2 baseline (0.9997). These results
demonstrate that Evo2 captures variant-effect information not redundant with
existing supervised annotations.

\subsection{Non-coding and indel variants: filling critical gaps}
\label{sec:results_noncoding}

A key advantage of genomic foundation models over protein-level tools is the
ability to score non-coding and indel variants using the same framework.

\subsubsection*{TERT promoter MPRA.}

We validated Evo2's non-coding predictions against massively parallel reporter
assay (MPRA) data for 777 variants in the TERT promoter
(\Cref{fig:noncoding}A). Evo2 $\Delta$LL correlated positively with MPRA
reporter activity (Spearman $\rho = +0.267$, $p = 3.8 \times 10^{-14}$),
indicating that variants predicted as more deleterious by Evo2 also reduced
transcriptional output in the GBM reporter assay. Critically, Evo2 was the
\textit{only} tool with a significant positive MPRA correlation: CADD showed
a non-significant negative trend ($\rho = -0.081$, $p = 0.21$, $n = 240$),
SpliceAI was near zero ($\rho = +0.003$, $p = 0.96$), and ncER showed a
significant \textit{negative} correlation ($\rho = -0.211$, $p = 5.9 \times 10^{-4}$).
These results highlight a genuine advantage of sequence-level foundation models
for non-coding variant interpretation.

\subsubsection*{ENCODE cCRE enrichment.}

Variants overlapping ENCODE candidate cis-regulatory elements showed elevated
Evo2 scores relative to non-cCRE variants. Enrichment ratios ranged from 1.12$\times$
(pELS) to 1.26$\times$ (dELS) and 1.19$\times$ (DNase-H3K4me3), consistent
with the expectation that regulatory elements harbour more functionally
consequential sequence.

\subsubsection*{Frameshift indel classification.}

Evo2 discriminated pathogenic from benign frameshift indels with high accuracy
across six genes with sufficient ClinVar-annotated indels
(\Cref{fig:noncoding}C; \Cref{fig:supp_indels}). Frameshift AUROCs ranged from 0.948 (TERT, $n = 108$)
to 0.993 (ATM, $n = 2{,}699$), with a mean of 0.977. This demonstrates that the
$\Delta$LL framework, which captures disruption to the sequence likelihood
landscape, naturally extends to length-altering variants without requiring
specialised indel models.

\subsection{Mutation signature analysis reveals transversion bias}
\label{sec:results_radiation}

To assess whether Evo2 scores reflect mutational signatures relevant to
spaceflight radiation exposure, we classified variants by mutational category
(\Cref{fig:supp_signatures}A). Across all ten genes, transversions scored
significantly more deleterious than transitions ($p < 10^{-10}$ for every gene),
consistent with the greater chemical disruption caused by purine$\leftrightarrow$pyrimidine
substitutions.

\subsubsection*{Transversion control.}

To test whether radiation-specific mutations (C$>$A and G$>$T, associated with
8-oxoguanine lesions from ionising radiation) score differently from other
transversions, we performed a within-transversion comparison
(\Cref{fig:supp_signatures}B; Supplementary Table~S6). After Bonferroni
correction for ten genes, only two showed marginal radiation-specific signal:
TP53 ($p_\text{Bonf} = 0.005$, rank-biserial $r_{bc} = +0.049$) and DNMT3A
($p_\text{Bonf} = 0.032$, $r_{bc} = +0.033$). The remaining eight genes showed
no significant difference between radiation-oxidative and other transversions
($p_\text{Bonf} > 0.32$). This control demonstrates that the consistently
elevated scores for radiation-type mutations reported in the four-class analysis
are primarily attributable to transversion-versus-transition bias rather than
radiation-specific chemistry.

\subsection{Variant atlas and constraint landscapes}
\label{sec:results_atlas}

In total, we scored 215,001 variants across the ten gene loci (168,409 SNVs and
46,592 indels), creating a comprehensive variant effect atlas for spaceflight
radiation-response genes. Gene-level constraint landscapes (\Cref{fig:landscapes})
reveal that functional protein domains correspond to regions of elevated predicted
deleteriousness, consistent with structural and biochemical expectations.

To illustrate clinical utility, we scored eight variants previously detected in
astronaut genomes from the Rutter et al.\ cohort~\cite{Rutter_2024}
(\Cref{fig:supp_astronauts}; Supplementary Table~S7). The somatic hotspot DNMT3A
R882C ranked at the 97.2nd percentile of all DNMT3A variants, and TP53 R248W
ranked at the 86.5th percentile. The TERT C228T promoter mutation ranked at only
the 4th percentile, consistent with its gain-of-function mechanism---this variant
\textit{increases} TERT expression by creating a de novo ETS binding site, and a
model trained to detect sequence disruption would not flag a variant that
preserves or enhances regulatory grammar.
